\begin{table}[h!]
\centering
% \caption{Managerial Synthesis of Research Implications for Resilience}
\label{tab:Implications}
\begin{tabular}{p{3cm} p{5cm} p{5cm}}
\hline
\textbf{Implication} & \textbf{Research Focus} & \textbf{Managerial Relevance} \\
\hline
\textbf{Research Direction} & \textbf{Scholarly Focus} & \textbf{Managerial Relevance} \\
\hline
Ontological shift: Resilience as emergent capability 
& Model organizations as adaptive socio-technical systems; theorize how routines, infrastructures, and institutional rules interact. 
& Clarifies that resilience is not robustness or redundancy, but an adaptive capability managers must deliberately cultivate. \\

Indicators and modeling: From metaphor to measurement 
& Develop behavioral and computational markers (e.g., coordination speed, reconfiguration capacity, variability under stress); explore via modeling and simulation. 
& Provides measurable benchmarks to assess preparedness, avoid vague rhetoric, and compare resilience across units or firms. \\

Empirical grounding: Linking theory to practice 
& Use longitudinal studies, digital trace data, and fieldwork to capture how resilience evolves and how institutional frictions constrain it. 
& Ensures metrics reflect lived organizational realities, enabling managers to validate frameworks against actual disruptions. \\

Integration across levels and methods 
& Design multi-method programs that connect micro routines, meso infrastructures, macro institutions, and inter-organizational ecosystems. 
& Helps managers see cross-level dependencies, avoiding siloed investments and fostering systemic resilience strategies. \\

Managerial design principles and actionable tools 
& Translate insights into benchmarks, early-warning systems, and governance mechanisms; prevent “metrics theater.” 
& Offers managers concrete instruments—dashboards, simulations, collaborative exercises—to build genuine, testable resilience. \\
\hline
\end{tabular}
\caption{\textbf{Managerial synthesis of research directions for advancing organizational resilience.} 
The table links five core directions identified in this review—ontological clarification, operational indicators, empirical grounding, cross-level integration, and managerial design principles—highlighting their relevance for both IS research and organizational practice.}
\label{tab:Implications}
\end{table}